\section{Глава~\ref{sec:sensors}. Входы машины}

Устройство датчиков измерено прямо --- записями тока одиночных клеток, колебаний улитки и подсчётом клеток сетчатки; предел чувствительности и формула дробового шума следуют из физики, а то, что коды датчиков подстроены под наибольшую информацию, --- гипотеза с сильной опорой.

{\footnotesize
\setlength{\tabcolsep}{3pt}\renewcommand{\arraystretch}{1.15}
\begin{longtable}{>{\raggedright}p{0.5cm}>{\raggedright}p{4.0cm}>{\raggedright}p{5.6cm}>{\raggedright}p{2.3cm}>{\raggedright\arraybackslash}p{1.7cm}}
\toprule
№ & Утверждение & Опыт и что измерено & Источник & Статус \\
\midrule\endfirsthead
\toprule
№ & Утверждение & Опыт и что измерено & Источник & Статус \\
\midrule\endhead
1 & Датчик --- преобразователь: рецептор даёт ток, а выход чувствительных клеток глаза и уха --- глутамат на шину \nt{ГЛУТ}; первые клетки импульсов не выдают (рис.~\ref{fig:transducer}) & Палочки и колбочки черепахи выбрасывают возбуждающую аминокислоту: её ловят глутаматные каналы на отсечённом кусочке мембраны. Внутренние волосковые клетки улитки крысы: токи в окончаниях нервных волокон идут через глутаматные AMPA-рецепторы, частота выбросов растёт с плавной деполяризацией клетки до $150$\,Гц & \cite{CopenhagenJahr1989,GlowatzkiFuchs2002} & измерено \\
2 & Датчик света в темноте отвечает на один фотон током около пикоампера & Всасывающий электрод на наружном сегменте палочки жабы: ответы на тусклые вспышки квантованы, событие $\approx1$\,пА ($5\%$ темнового тока) с разбросом $0.2$\,пА, эффективность $0.5$. Человек сообщает о попадании одного фотона чаще случайного & \cite{Baylor1979,Tinsley2016} & измерено \\
3 & Датчик звука замечает смещение меньше нанометра & Метод Мёссбауэра, основная мембрана улитки морской свинки на месте $16$--$19$\,кГц: на пороге ответа слухового нерва скорость мембраны около $0.04$\,мм/с, то есть смещение $\approx0.35$\,нм (пересчёт наш). Мерили мембрану, а не пучок датчика & \cite{Sellick1982,Hudspeth1989} & измерено \\
4 & Точность в темноте ограничена дробовым шумом фотонов~\eqref{eq:photonnoise}; при слабом свете накопление длиннее & Формула следует из пуассоновской статистики. В палочке шум при тусклом свете совпадает с суммой независимых квантовых событий; на ярком фоне ответ на вспышку меньше и короче. Число «десятые доли секунды» здесь отдельно не проверено & \cite{Baylor1979} & теория \\
5 & Диапазон входа --- десять порядков для света и двенадцать для звука, а частота импульсов --- меньше трёх порядков & Для звука: ответ основной мембраны на характеристической частоте растёт с крутизной до $0.2$\,дБ/дБ в полосе $40$--$80$\,дБ, сжатие видно и выше $100$\,дБ. Сами числа порядков --- стандартные оценки, в главе без источника & \cite{Ruggero1997} & измерено \\
6 & Ответ датчика~\eqref{eq:nakarushton}, а $\sigma$ следует за средней яркостью, сдвигая кривую по логарифмической оси (рис.~\ref{fig:adaptation}) & Формула впервые подобрана к S-потенциалам сетчатки рыбы. Колбочки черепахи: ответы подчиняются $V=I V_m/(I+\sigma)$ и охватывают $\sim3.5$ порядка яркости; после $2$--$3$ минут фона кривая сдвигается по горизонтали, сохраняя ширину. Связь с делением на общую активность --- обзор & \cite{NakaRushton1966,NormannPerlman1979,Carandini2012} & измерено \\
7 & Датчики передают перемены, а их коды подстроены под наибольшую информацию на импульс (утверждение~\ref{prop:sensors}) & Принцип предложен теоретически. Сильнейшее свидетельство: у мухи характеристика «контраст $\to$ ответ» нейронов первого слоя совпадает с накопленным распределением контрастов природных сцен --- так достигается наибольшая информация & \cite{Barlow1961,Laughlin1981} & гипотеза \\
8 & Включающие и выключающие датчики --- двухпроводный код; событийная камера повторяет его в кремнии & Обзор опытов: включающий и выключающий каналы сетчатки разделяются уже на первом синапсе и выключаются по отдельности. Камера: $128\times128$ пикселей без кадров, диапазон $120$\,дБ, задержка $15$\,мкс & \cite{Schiller1992,Lichtsteiner2008,Gallego2022} & измерено \\
9 & В глазу $\sim10^8$ датчиков света и $\sim10^6$ выходных нейронов: сжатие в $\sim100$ раз & Подсчёт на целых препаратах сетчатки человека: в среднем $92$ млн палочек и $4.6$ млн колбочек; выходных (ганглиозных) клеток от $0.7$ до $1.5$ млн в разных глазах & \cite{Curcio1990,CurcioAllen1990} & измерено \\
10 & У мыши выходных каналов глаза больше тридцати & Мышь: двухфотонная кальциевая съёмка более $11\,000$ выходных клеток и кластеризация ответов --- заметно больше $30$ функциональных типов. Для человека число не измерено & \cite{Baden2016} & измерено \\
11 & Глаз морской свинки передаёт $\sim875\,000$ бит/с; пересчёт на глаз человека даёт $\sim10^7$ бит/с & Морская свинка, многоэлектродная матрица, движение в природных сценах: $6$--$13$ бит/с на клетку, $\sim875\,000$ бит/с на $\sim10^5$ выходных клеток. Для человека ($\sim10^6$ клеток) $\sim10^7$ бит/с --- пересчёт авторов & \cite{Koch2006} & измерено \\
12 & Ухо --- банк полосовых фильтров с почти логарифмической картой частот (рис.~\ref{fig:filterbank}) & Место наибольшего колебания основной мембраны зависит от частоты; функция «частота--место» почти экспоненциальна, одной формой описывает улитки человека и живых животных & \cite{Greenwood1990,Robles2001} & измерено \\
13 & Резонаторы активны: часть датчиков усиливает слабые колебания в тысячи раз по амплитуде ($66$--$76$\,дБ), с ростом громкости усиление падает & Лазерная виброметрия основания улитки шиншиллы: при слабых тонах на характеристической частоте усиление $66$--$76$\,дБ относительно стремечка; смерть снижает чувствительность на $60$--$81$\,дБ (в $10^3$--$10^4$ раз по амплитуде) и убирает сжатие. Источник силы --- наружные волосковые клетки & \cite{Ruggero1997,Robles2001} & измерено \\
14 & Усилитель уха стоит у границы самовозбуждения & Расчёт: у точки бифуркации Хопфа неизбежны сжатие диапазона, острая настройка и комбинационные тоны --- все известные нелинейности слуха. Что улитка настроена именно так, прямо не доказано & \cite{Eguiluz2000} & гипотеза \\
15 & На низких частотах импульсы идут в фазе со звуком (у морской свинки привязка слабеет от $\sim600$\,Гц и заметна до $\sim3.5$\,кГц), и по ним находится направление; на высоких остаётся код местом & Слуховой нерв морской свинки: привязка к фазе начинает падать около $600$\,Гц и исчезает выше $3.5$\,кГц; у кошки и обезьяны граница примерно на октаву выше. Разница времени прихода в два уха --- обзор & \cite{PalmerRussell1986,Grothe2010} & измерено \\
16 & Остальные датчики (табл.~\ref{tab:sensors}): у обоняния $\sim400$ типов рецепторов, по одному на нейрон, и комбинаторный код; вкус частично через АТФ и серотонин; осязание --- быстро и медленно привыкающие датчики; тепло и холод --- отдельные & Мышь: один рецептор узнаёт много запахов, один запах --- много рецепторов, разные запахи --- разные наборы. Без рецепторов АТФ нерв вкуса не отвечает; вкусовые почки выбрасывают серотонин. Запись одиночных волокон кожи руки человека: четыре типа по скорости привыкания. Холодовой канал отличен от тепловых & \cite{Mombaerts2004,Malnic1999,Finger2005,Huang2005,JohanssonVallbo1979,McKemy2002} & измерено \\
\bottomrule
\end{longtable}}
